% GENERATED FILE. Edit canonical lessons, then run npm run generate:books.
% canonical-source-hash: fnv1a64:16104cc0d22aa766
% canonical-lessons: GU-C290-daku, GU-C290-banduk, GU-C290-dhamki, GU-C290-dhamkavvu, GU-C290-chinavvu

\chapter{Bandits on the Road}
\label{ch:gu-bandits-on-the-road}
\hlchaptermodality{290}

\begin{chapteropening}
\textbf{By the end of this chapter:} \emph{I can tell how bandits stopped a bus and what they took.}

Complete the last lesson of chapter 290: I can tell how bandits stopped a bus and what they took.
\end{chapteropening}

\section[\texorpdfstring{\gu{ડાકુ} (ḍāku)}{ડાકુ (ḍāku)}]{\gu{ડાકુ} (ḍāku) — a bandit}

\label{lesson:GU-C290-daku}



\begin{quote}
Before the new one: say the Gujarati for a punishment, then the Gujarati for to undergo.
\end{quote}

\subsection*{You'll want to know: \gu{ડાકુ}}
\textbf{\gu{ડાકુ}} — \emph{ḍāku} — \textquotedblleft{}a bandit\textquotedblright{}.

\begin{quote}
\gu{જૂના} \gu{વખતમાં} \gu{આ} \gu{જંગલમાં} \gu{ડાકુ} \gu{રહેતા} \gu{હતા}.
\end{quote}

\begin{quote}
\emph{jūnā vakhatmā̃ ā jangalmā̃ ḍāku rahetā hatā.}
\end{quote}

\begin{quote}
In the old days bandits lived in this forest.
\end{quote}

A masculine noun.

\begin{cousinweb}
comes from \emph{*ḍākka}, \textquotedblleft{}robber, robbery\textquotedblright{}, a form CDIAL reconstructs; CDIAL lists Gujarati \emph{ḍāko}, an attack by robbers, in the same entry (CDIAL 5543). A \textbf{\gu{ડાકુ}} belongs to a gang that robs on the roads.
\end{cousinweb}

\begin{grammarlens}[title={What you've built}]
One more word for telling what happened.
\end{grammarlens}

\subsection*{Your turn}
\begin{itemize}
\raggedright
  \item \emph{Say it:} \emph{ḍāku}
  \item \emph{Say it:} \emph{ḍāku}, once more
  \item \emph{Say it:} \emph{jūnā vakhatmā̃ ā jangalmā̃ ḍāku rahetā hatā.} — In the old days bandits lived in this forest.
  \item \emph{Recall:} say the Gujarati for a punishment, then the Gujarati for to undergo, then say \emph{ḍāku} again
\end{itemize}

\begin{tcolorbox}[breakable,colback=teal!4,colframe=teal!35!black,title={Before you move on}]
What does \gu{ડાકુ} mean? (\textquotedblleft{}A bandit\textquotedblright{}.) Say it once more.
\end{tcolorbox}
\section[\texorpdfstring{\gu{બંદૂક} (bandūk)}{બંદૂક (bandūk)}]{\gu{બંદૂક} (bandūk) — a gun}

\label{lesson:GU-C290-banduk}



\begin{quote}
Before the new one: say the Gujarati for to undergo, then the Gujarati for a bandit.
\end{quote}

\subsection*{You'll want to know: \gu{બંદૂક}}
\textbf{\gu{બંદૂક}} — \emph{bandūk} — \textquotedblleft{}a gun\textquotedblright{}.

\begin{quote}
\gu{ડાકુએ} \gu{કહ્યું} \gu{કે} \gu{જો} \gu{કોઈ} \gu{ભાગશે} \gu{તો} \gu{તે} \gu{બંદૂક} \gu{ચલાવશે}.
\end{quote}

\begin{quote}
\emph{ḍākue kahyũ ke jo koī bhāgśe to te bandūk calāvśe.}
\end{quote}

\begin{quote}
The bandit said that if anyone ran, he would fire his gun.
\end{quote}

A feminine noun.

\begin{cousinweb}
\textbf{\gu{બંદૂક}} (\emph{bandūk}) is usually traced, through Persian, to Arabic \emph{bunduq}, a hazelnut and then a small round pellet; the step to a gun is said to have come through the pellet bow. It is feminine, and a gun is fired with \textbf{\gu{ચલાવવું}}: \textbf{\gu{બંદૂક} \gu{ચલાવી}}.
\end{cousinweb}

\begin{grammarlens}[title={What you've built}]
One more word for telling what happened.
\end{grammarlens}

\subsection*{Your turn}
\begin{itemize}
\raggedright
  \item \emph{Say it:} \emph{bandūk}
  \item \emph{Say it:} \emph{bandūk}, once more
  \item \emph{Say it:} \emph{ḍākue kahyũ ke jo koī bhāgśe to te bandūk calāvśe.} — The bandit said that if anyone ran, he would fire his gun.
  \item \emph{Recall:} say the Gujarati for to undergo, then the Gujarati for a bandit, then say \emph{bandūk} again
\end{itemize}

\begin{tcolorbox}[breakable,colback=teal!4,colframe=teal!35!black,title={Before you move on}]
What does \gu{બંદૂક} mean? (\textquotedblleft{}A gun\textquotedblright{}.) Say it once more.
\end{tcolorbox}
\section[\texorpdfstring{\gu{ધમકી} (dhamkī)}{ધમકી (dhamkī)}]{\gu{ધમકી} (dhamkī) — a threat}

\label{lesson:GU-C290-dhamki}



\begin{quote}
Before the new one: say the Gujarati for a bandit, then the Gujarati for a gun.
\end{quote}

\subsection*{You'll want to know: \gu{ધમકી}}
\textbf{\gu{ધમકી}} — \emph{dhamkī} — \textquotedblleft{}a threat\textquotedblright{}.

\begin{quote}
\gu{ડાકુઓએ} \gu{ડ્રાઈવરને} \gu{ધમકી} \gu{આપી}.
\end{quote}

\begin{quote}
\emph{ḍākuoe ḍrāīvarne dhamkī āpī.}
\end{quote}

\begin{quote}
The bandits threatened the driver.
\end{quote}

A feminine noun.

\begin{cousinweb}
\textbf{\gu{ધમકી}} (\emph{dhamkī}) comes from \emph{*dhammakka}, \textquotedblleft{}a threat\textquotedblright{}, a form CDIAL reconstructs, and CDIAL lists the Gujarati noun (CDIAL 6736). You give a threat: \textbf{\gu{ધમકી} \gu{આપી}}.
\end{cousinweb}

\begin{grammarlens}[title={What you've built}]
One more word for telling what happened.
\end{grammarlens}

\subsection*{Your turn}
\begin{itemize}
\raggedright
  \item \emph{Say it:} \emph{dhamkī}
  \item \emph{Say it:} \emph{dhamkī}, once more
  \item \emph{Say it:} \emph{ḍākuoe ḍrāīvarne dhamkī āpī.} — The bandits threatened the driver.
  \item \emph{Recall:} say the Gujarati for a bandit, then the Gujarati for a gun, then say \emph{dhamkī} again
\end{itemize}

\begin{tcolorbox}[breakable,colback=teal!4,colframe=teal!35!black,title={Before you move on}]
What does \gu{ધમકી} mean? (\textquotedblleft{}A threat\textquotedblright{}.) Say it once more.
\end{tcolorbox}
\section[\texorpdfstring{\gu{ધમકાવવું} (dhamkāvvũ)}{ધમકાવવું (dhamkāvvũ)}]{\gu{ધમકાવવું} (dhamkāvvũ) — to threaten}

\label{lesson:GU-C290-dhamkavvu}



\begin{quote}
Before the new one: say the Gujarati for a gun, then the Gujarati for a threat.
\end{quote}

\subsection*{You'll want to know: \gu{ધમકાવવું}}
\textbf{\gu{ધમકાવવું}} — \emph{dhamkāvvũ} — \textquotedblleft{}to threaten\textquotedblright{}.

\begin{quote}
\gu{એક} \gu{ડાકુએ} \gu{બંદૂક} \gu{બતાવીને} \gu{મુસાફરોને} \gu{ધમકાવ્યા}.
\end{quote}

\begin{quote}
\emph{ek ḍākue bandūk batāvīne musāphrone dhamkāvyā.}
\end{quote}

\begin{quote}
One bandit showed his gun and threatened the passengers.
\end{quote}

Past, with \gu{તેણે} (he or she …): \gu{તેણે} \gu{મુસાફરોને} \gu{ધમકાવ્યા}, \emph{teṇe musāphrone dhamkāvyā}.

\begin{cousinweb}
\textbf{\gu{ધમકાવવું}} (\emph{dhamkāvvũ}) belongs with \textbf{\gu{ધમકી}}, a threat: CDIAL lists the Gujarati verb in the same entry, \emph{*dhammakka} (CDIAL 6736). To threaten is to frighten someone with words or a weapon.
\end{cousinweb}

\begin{grammarlens}[title={What you've built}]
One more word for telling what happened.
\end{grammarlens}

\subsection*{Your turn}
\begin{itemize}
\raggedright
  \item \emph{Say it:} \emph{dhamkāvvũ}
  \item \emph{Say it:} \emph{dhamkāvvũ}, once more
  \item \emph{Say it:} \emph{ek ḍākue bandūk batāvīne musāphrone dhamkāvyā.} — One bandit showed his gun and threatened the passengers.
  \item \emph{Recall:} say the Gujarati for a gun, then the Gujarati for a threat, then say \emph{dhamkāvvũ} again
\end{itemize}

\begin{tcolorbox}[breakable,colback=teal!4,colframe=teal!35!black,title={Before you move on}]
What does \gu{ધમકાવવું} mean? (\textquotedblleft{}To threaten\textquotedblright{}.) Say it once more.
\end{tcolorbox}
\section[\texorpdfstring{\gu{છીનવવું} (chīnavvũ)}{છીનવવું (chīnavvũ)}]{\gu{છીનવવું} (chīnavvũ) — to snatch}

\label{lesson:GU-C290-chinavvu}



\begin{quote}
Before the new one: say the Gujarati for a threat, then the Gujarati for to threaten.
\end{quote}

\subsection*{You'll want to know: \gu{છીનવવું}}
\textbf{\gu{છીનવવું}} — \emph{chīnavvũ} — \textquotedblleft{}to snatch\textquotedblright{}.

\begin{quote}
\gu{ડાકુઓએ} \gu{મારી} \gu{થેલી} \gu{છીનવી} \gu{લીધી}.
\end{quote}

\begin{quote}
\emph{ḍākuoe mārī thelī chīnvī līdhī.}
\end{quote}

\begin{quote}
The bandits snatched my bag.
\end{quote}

Past, with \gu{તેણે} (he or she …): \gu{તેણે} \gu{થેલી} \gu{છીનવી}, \emph{teṇe thelī chīnvī}.

\begin{cousinweb}
\textbf{\gu{છીનવવું}} (\emph{chīnavvũ}) is set by CDIAL under Sanskrit \emph{chinná}, \textquotedblleft{}cut off\textquotedblright{}, where it lists the Gujarati verb (CDIAL 5047): what is snatched is torn away from its owner. It often takes \textbf{\gu{લેવું}}: \textbf{\gu{છીનવી} \gu{લીધી}}, snatched it away.
\end{cousinweb}

\begin{grammarlens}[title={What you've built}]
One more word for telling what happened.
\end{grammarlens}

\subsection*{Your turn}
\begin{itemize}
\raggedright
  \item \emph{Say it:} \emph{chīnavvũ}
  \item \emph{Say it:} \emph{chīnavvũ}, once more
  \item \emph{Say it:} \emph{ḍākuoe mārī thelī chīnvī līdhī.} — The bandits snatched my bag.
  \item \emph{Recall:} say the Gujarati for a threat, then the Gujarati for to threaten, then say \emph{chīnavvũ} again
\end{itemize}

\begin{tcolorbox}[breakable,colback=teal!4,colframe=teal!35!black,title={Before you move on}]
What does \gu{છીનવવું} mean? (\textquotedblleft{}To snatch\textquotedblright{}.) Say it once more.
\end{tcolorbox}
